% chapter_skeleton.tex -- copy to <class>/<chapter>/latex/main.tex and fill in.
% Copy ../skill/templates/preamble.tex next to it (or adjust the \input path).
\documentclass[11pt,a4paper]{article}
% preamble.tex -- shared preamble for NCERT board-style solution books.
% Use:  \documentclass[11pt,a4paper]{article}  % preamble.tex -- shared preamble for NCERT board-style solution books.
% Use:  \documentclass[11pt,a4paper]{article}  % preamble.tex -- shared preamble for NCERT board-style solution books.
% Use:  \documentclass[11pt,a4paper]{article}  \input{preamble}  \begin{document} ... \end{document}
% Compile with pdflatex (run twice: longtable and cross-references need a second pass).

\usepackage{amsmath,amssymb,graphicx,enumitem,tikz,xcolor,tabularx,array,needspace,longtable}
\usetikzlibrary{calc,angles,quotes}
\usepackage[margin=2.2cm]{geometry}
\setlength{\parindent}{0pt}
\setlength{\parskip}{5pt}
\renewcommand{\arraystretch}{1.2}

% ---- shorthand for geometry write-ups ----
\renewcommand{\deg}{^\circ}          % \deg already exists in LaTeX, so \renewcommand
\newcommand{\ang}{\angle}
\newcommand{\tri}{\triangle}

% ---- never leave a question heading stranded at the foot of a page ----
\let\oldsubsection\subsection
\renewcommand{\subsection}{\Needspace{12\baselineskip}\oldsubsection}

% ---- "Proof:" heading that stays with the table that follows ----
\newcommand{\proofhd}{\par\Needspace{6\baselineskip}\textbf{Proof:}}

% ---- board-style proof table: Step | Statement | Reason ----
% Usage:  \begin{prf} 1 & statement & reason \\ 2 & ... \end{prf}      (no trailing \\)
% longtable, so a long proof may break across pages (the header repeats).
\setlength{\LTpre}{2pt}\setlength{\LTpost}{4pt}
\newenvironment{prf}
  {\par\noindent\renewcommand{\arraystretch}{1.4}%
   \longtable{@{}c|>{\raggedright\arraybackslash}p{0.46\linewidth}|>{\raggedright\arraybackslash}p{0.38\linewidth}@{}}%
   \hline \textbf{Step} & \textbf{Statement} & \textbf{Reason}\\ \hline \endfirsthead
   \hline \textbf{Step} & \textbf{Statement} & \textbf{Reason (continued)}\\ \hline \endhead}
  {\\ \hline\endlongtable}

% ---- figure helpers ----
% \dotat{P}{above}  : a dot at coordinate (P) with the label placed above it
% \figcap{text}     : caption line under a figure (the \\* keeps it on the same page)
\newcommand{\dotat}[2]{\fill (#1) circle (1.6pt) node[#2]{\small $#1$};}
\newcommand{\figcap}[1]{\\*[-2pt]{\small #1}}

% ---- final result of a solution, highlighted at the END of the Solution ----
% Solution and answer are ONE block: write the working, then \ans{...}. No separate "Answer." heading.
\newcommand{\ans}[1]{\par\vspace{4pt}\noindent\fcolorbox{black!55}{blue!6}{\parbox{\dimexpr\linewidth-2\fboxsep-2\fboxrule\relax}{\raggedright\textbf{Answer:} #1}}\par\vspace{2pt}}

% ---- common-mistake note for the student who went wrong at this step ----
% Usage: \slip{what is usually done wrongly, and why it is wrong}
\newcommand{\slip}[1]{\par\noindent{\small\itshape\textbf{Common mistake.} #1}\par}

% ---- distance formula written out in full, so every substituted value can be checked ----
% \bk{n}  : n, in brackets when n is negative
% \dist{AB}{x1}{y1}{x2}{y2}{simplified value}  : one row group of an align* (add \\ after it)
% \distn{AB}{x1}{y1}{x2}{y2}                   : same, but stops at the unsimplified root
% Rows: (1) substitute into the formula, (2) work out each bracket,
%       (3) square each difference, (4) add the squares and take the root.
% The numbers are computed by TeX (\numexpr), so the displayed values cannot disagree with the formula.
% Use \distnote once before the first block to tell the student how to read the rows.
\newcommand{\bk}[1]{\ifnum#1<0 (#1)\else #1\fi}
\newcommand{\distcore}[5]{#1&=\sqrt{(#4-\bk{#2})^2+(#5-\bk{#3})^2}\\
&=\sqrt{\bk{\the\numexpr#4-#2\relax}^2+\bk{\the\numexpr#5-#3\relax}^2}\\
&=\sqrt{\the\numexpr(#4-#2)*(#4-#2)\relax+\the\numexpr(#5-#3)*(#5-#3)\relax}
=\sqrt{\the\numexpr(#4-#2)*(#4-#2)+(#5-#3)*(#5-#3)\relax}}
\newcommand{\distnote}{\textit{How to read the working below:} in each pair the first point named is $(x_1,y_1)$ and the second is $(x_2,y_2)$. Line 1 puts the values into the formula, line 2 works out each bracket, line 3 squares the brackets, and then the squares are added and the square root is taken.}
\newcommand{\dist}[6]{\distcore{#1}{#2}{#3}{#4}{#5}=#6}
\newcommand{\distn}[5]{\distcore{#1}{#2}{#3}{#4}{#5}}

% ---- NCERT-style 10 x 10 grid with Columns along the bottom and Rows up the side ----
% \ncertgrid{ <tikz commands drawing only the GIVEN points> }  -- never print solved coordinates
\newcommand{\ncertgrid}[1]{%
\begin{tikzpicture}[scale=0.5]
\fill[cyan!12] (0,0) rectangle (10,10);
\draw[step=1,gray!60,very thin] (0,0) grid (10,10);
\draw[thick] (0,0) rectangle (10,10);
\foreach \i in {0,...,10}{\node[below,font=\scriptsize] at (\i,0) {\i};}
\foreach \i in {1,...,10}{\node[left,font=\scriptsize] at (0,\i) {\i};}
\node[below,font=\small\bfseries] at (5,-0.8) {Columns};
\node[rotate=90,above,font=\small\bfseries] at (-1.0,5) {Rows};
#1
\end{tikzpicture}}
  \begin{document} ... \end{document}
% Compile with pdflatex (run twice: longtable and cross-references need a second pass).

\usepackage{amsmath,amssymb,graphicx,enumitem,tikz,xcolor,tabularx,array,needspace,longtable}
\usetikzlibrary{calc,angles,quotes}
\usepackage[margin=2.2cm]{geometry}
\setlength{\parindent}{0pt}
\setlength{\parskip}{5pt}
\renewcommand{\arraystretch}{1.2}

% ---- shorthand for geometry write-ups ----
\renewcommand{\deg}{^\circ}          % \deg already exists in LaTeX, so \renewcommand
\newcommand{\ang}{\angle}
\newcommand{\tri}{\triangle}

% ---- never leave a question heading stranded at the foot of a page ----
\let\oldsubsection\subsection
\renewcommand{\subsection}{\Needspace{12\baselineskip}\oldsubsection}

% ---- "Proof:" heading that stays with the table that follows ----
\newcommand{\proofhd}{\par\Needspace{6\baselineskip}\textbf{Proof:}}

% ---- board-style proof table: Step | Statement | Reason ----
% Usage:  \begin{prf} 1 & statement & reason \\ 2 & ... \end{prf}      (no trailing \\)
% longtable, so a long proof may break across pages (the header repeats).
\setlength{\LTpre}{2pt}\setlength{\LTpost}{4pt}
\newenvironment{prf}
  {\par\noindent\renewcommand{\arraystretch}{1.4}%
   \longtable{@{}c|>{\raggedright\arraybackslash}p{0.46\linewidth}|>{\raggedright\arraybackslash}p{0.38\linewidth}@{}}%
   \hline \textbf{Step} & \textbf{Statement} & \textbf{Reason}\\ \hline \endfirsthead
   \hline \textbf{Step} & \textbf{Statement} & \textbf{Reason (continued)}\\ \hline \endhead}
  {\\ \hline\endlongtable}

% ---- figure helpers ----
% \dotat{P}{above}  : a dot at coordinate (P) with the label placed above it
% \figcap{text}     : caption line under a figure (the \\* keeps it on the same page)
\newcommand{\dotat}[2]{\fill (#1) circle (1.6pt) node[#2]{\small $#1$};}
\newcommand{\figcap}[1]{\\*[-2pt]{\small #1}}

% ---- final result of a solution, highlighted at the END of the Solution ----
% Solution and answer are ONE block: write the working, then \ans{...}. No separate "Answer." heading.
\newcommand{\ans}[1]{\par\vspace{4pt}\noindent\fcolorbox{black!55}{blue!6}{\parbox{\dimexpr\linewidth-2\fboxsep-2\fboxrule\relax}{\raggedright\textbf{Answer:} #1}}\par\vspace{2pt}}

% ---- common-mistake note for the student who went wrong at this step ----
% Usage: \slip{what is usually done wrongly, and why it is wrong}
\newcommand{\slip}[1]{\par\noindent{\small\itshape\textbf{Common mistake.} #1}\par}

% ---- distance formula written out in full, so every substituted value can be checked ----
% \bk{n}  : n, in brackets when n is negative
% \dist{AB}{x1}{y1}{x2}{y2}{simplified value}  : one row group of an align* (add \\ after it)
% \distn{AB}{x1}{y1}{x2}{y2}                   : same, but stops at the unsimplified root
% Rows: (1) substitute into the formula, (2) work out each bracket,
%       (3) square each difference, (4) add the squares and take the root.
% The numbers are computed by TeX (\numexpr), so the displayed values cannot disagree with the formula.
% Use \distnote once before the first block to tell the student how to read the rows.
\newcommand{\bk}[1]{\ifnum#1<0 (#1)\else #1\fi}
\newcommand{\distcore}[5]{#1&=\sqrt{(#4-\bk{#2})^2+(#5-\bk{#3})^2}\\
&=\sqrt{\bk{\the\numexpr#4-#2\relax}^2+\bk{\the\numexpr#5-#3\relax}^2}\\
&=\sqrt{\the\numexpr(#4-#2)*(#4-#2)\relax+\the\numexpr(#5-#3)*(#5-#3)\relax}
=\sqrt{\the\numexpr(#4-#2)*(#4-#2)+(#5-#3)*(#5-#3)\relax}}
\newcommand{\distnote}{\textit{How to read the working below:} in each pair the first point named is $(x_1,y_1)$ and the second is $(x_2,y_2)$. Line 1 puts the values into the formula, line 2 works out each bracket, line 3 squares the brackets, and then the squares are added and the square root is taken.}
\newcommand{\dist}[6]{\distcore{#1}{#2}{#3}{#4}{#5}=#6}
\newcommand{\distn}[5]{\distcore{#1}{#2}{#3}{#4}{#5}}

% ---- NCERT-style 10 x 10 grid with Columns along the bottom and Rows up the side ----
% \ncertgrid{ <tikz commands drawing only the GIVEN points> }  -- never print solved coordinates
\newcommand{\ncertgrid}[1]{%
\begin{tikzpicture}[scale=0.5]
\fill[cyan!12] (0,0) rectangle (10,10);
\draw[step=1,gray!60,very thin] (0,0) grid (10,10);
\draw[thick] (0,0) rectangle (10,10);
\foreach \i in {0,...,10}{\node[below,font=\scriptsize] at (\i,0) {\i};}
\foreach \i in {1,...,10}{\node[left,font=\scriptsize] at (0,\i) {\i};}
\node[below,font=\small\bfseries] at (5,-0.8) {Columns};
\node[rotate=90,above,font=\small\bfseries] at (-1.0,5) {Rows};
#1
\end{tikzpicture}}
  \begin{document} ... \end{document}
% Compile with pdflatex (run twice: longtable and cross-references need a second pass).

\usepackage{amsmath,amssymb,graphicx,enumitem,tikz,xcolor,tabularx,array,needspace,longtable}
\usetikzlibrary{calc,angles,quotes}
\usepackage[margin=2.2cm]{geometry}
\setlength{\parindent}{0pt}
\setlength{\parskip}{5pt}
\renewcommand{\arraystretch}{1.2}

% ---- shorthand for geometry write-ups ----
\renewcommand{\deg}{^\circ}          % \deg already exists in LaTeX, so \renewcommand
\newcommand{\ang}{\angle}
\newcommand{\tri}{\triangle}

% ---- never leave a question heading stranded at the foot of a page ----
\let\oldsubsection\subsection
\renewcommand{\subsection}{\Needspace{12\baselineskip}\oldsubsection}

% ---- "Proof:" heading that stays with the table that follows ----
\newcommand{\proofhd}{\par\Needspace{6\baselineskip}\textbf{Proof:}}

% ---- board-style proof table: Step | Statement | Reason ----
% Usage:  \begin{prf} 1 & statement & reason \\ 2 & ... \end{prf}      (no trailing \\)
% longtable, so a long proof may break across pages (the header repeats).
\setlength{\LTpre}{2pt}\setlength{\LTpost}{4pt}
\newenvironment{prf}
  {\par\noindent\renewcommand{\arraystretch}{1.4}%
   \longtable{@{}c|>{\raggedright\arraybackslash}p{0.46\linewidth}|>{\raggedright\arraybackslash}p{0.38\linewidth}@{}}%
   \hline \textbf{Step} & \textbf{Statement} & \textbf{Reason}\\ \hline \endfirsthead
   \hline \textbf{Step} & \textbf{Statement} & \textbf{Reason (continued)}\\ \hline \endhead}
  {\\ \hline\endlongtable}

% ---- figure helpers ----
% \dotat{P}{above}  : a dot at coordinate (P) with the label placed above it
% \figcap{text}     : caption line under a figure (the \\* keeps it on the same page)
\newcommand{\dotat}[2]{\fill (#1) circle (1.6pt) node[#2]{\small $#1$};}
\newcommand{\figcap}[1]{\\*[-2pt]{\small #1}}

% ---- final result of a solution, highlighted at the END of the Solution ----
% Solution and answer are ONE block: write the working, then \ans{...}. No separate "Answer." heading.
\newcommand{\ans}[1]{\par\vspace{4pt}\noindent\fcolorbox{black!55}{blue!6}{\parbox{\dimexpr\linewidth-2\fboxsep-2\fboxrule\relax}{\raggedright\textbf{Answer:} #1}}\par\vspace{2pt}}

% ---- common-mistake note for the student who went wrong at this step ----
% Usage: \slip{what is usually done wrongly, and why it is wrong}
\newcommand{\slip}[1]{\par\noindent{\small\itshape\textbf{Common mistake.} #1}\par}

% ---- distance formula written out in full, so every substituted value can be checked ----
% \bk{n}  : n, in brackets when n is negative
% \dist{AB}{x1}{y1}{x2}{y2}{simplified value}  : one row group of an align* (add \\ after it)
% \distn{AB}{x1}{y1}{x2}{y2}                   : same, but stops at the unsimplified root
% Rows: (1) substitute into the formula, (2) work out each bracket,
%       (3) square each difference, (4) add the squares and take the root.
% The numbers are computed by TeX (\numexpr), so the displayed values cannot disagree with the formula.
% Use \distnote once before the first block to tell the student how to read the rows.
\newcommand{\bk}[1]{\ifnum#1<0 (#1)\else #1\fi}
\newcommand{\distcore}[5]{#1&=\sqrt{(#4-\bk{#2})^2+(#5-\bk{#3})^2}\\
&=\sqrt{\bk{\the\numexpr#4-#2\relax}^2+\bk{\the\numexpr#5-#3\relax}^2}\\
&=\sqrt{\the\numexpr(#4-#2)*(#4-#2)\relax+\the\numexpr(#5-#3)*(#5-#3)\relax}
=\sqrt{\the\numexpr(#4-#2)*(#4-#2)+(#5-#3)*(#5-#3)\relax}}
\newcommand{\distnote}{\textit{How to read the working below:} in each pair the first point named is $(x_1,y_1)$ and the second is $(x_2,y_2)$. Line 1 puts the values into the formula, line 2 works out each bracket, line 3 squares the brackets, and then the squares are added and the square root is taken.}
\newcommand{\dist}[6]{\distcore{#1}{#2}{#3}{#4}{#5}=#6}
\newcommand{\distn}[5]{\distcore{#1}{#2}{#3}{#4}{#5}}

% ---- NCERT-style 10 x 10 grid with Columns along the bottom and Rows up the side ----
% \ncertgrid{ <tikz commands drawing only the GIVEN points> }  -- never print solved coordinates
\newcommand{\ncertgrid}[1]{%
\begin{tikzpicture}[scale=0.5]
\fill[cyan!12] (0,0) rectangle (10,10);
\draw[step=1,gray!60,very thin] (0,0) grid (10,10);
\draw[thick] (0,0) rectangle (10,10);
\foreach \i in {0,...,10}{\node[below,font=\scriptsize] at (\i,0) {\i};}
\foreach \i in {1,...,10}{\node[left,font=\scriptsize] at (0,\i) {\i};}
\node[below,font=\small\bfseries] at (5,-0.8) {Columns};
\node[rotate=90,above,font=\small\bfseries] at (-1.0,5) {Rows};
#1
\end{tikzpicture}}


\begin{document}

\begin{center}
{\Large\bfseries Class 10 Mathematics (NCERT) -- Chapter N: Title}\\[4pt]
{\large Worked solutions: Examples a--b, Exercise N.1 and Exercise N.2}\\[2pt]
{\small Draft for teacher verification.}
\end{center}

\subsection*{Formulas / properties used in this chapter}
\begin{enumerate}[label=(F\arabic*),leftmargin=*]
\item \textbf{Name of formula.} State it in general form, with the meaning of every symbol.
\end{enumerate}

% ------------------------------------------------------------------
% Pattern A: calculation / MCQ / fill-in
% Solution and answer are ONE block. The highlighted \ans box closes it.
% ------------------------------------------------------------------
\subsection*{Exercise N.1, Question 1}
\textbf{Question.} (exact NCERT wording, with options (A)--(D) if any)

% optional figure -- givens only, never the answer
% \begin{center}\begin{tikzpicture} ... \end{tikzpicture}%
% \figcap{Caption.}\end{center}

\textbf{Solution.} Name the formula in general form, then substitute and simplify.
Show every value BEFORE it is squared, rooted or multiplied, one step per line:
\begin{align*}
PQ^2&=OQ^2-OP^2\\
&=12^2-5^2\\
&=144-25=119
\end{align*}
\slip{The usual wrong turn at this step, and why it is wrong.}

\textbf{Verification.} Put the result back into the formula and show it agrees.

\ans{Option (X): value with unit.}

% ------------------------------------------------------------------
% Pattern B: proof
% ------------------------------------------------------------------
\subsection*{Exercise N.2, Question 2}
\textbf{Question.} Prove that \dots

\textbf{Solution.}

\textbf{Given:} \dots

\textbf{To prove:} \dots

\textbf{Construction:} \dots   % omit if none

\proofhd{}
\begin{prf}
1 & statement & reason (name the property in words) \\
2 & statement & reason
\end{prf}
\slip{The step students most often get wrong in this proof.}

\ans{The result stated in full. (Proved)}

% ------------------------------------------------------------------
% Pattern C: distance-formula style numeric blocks (coordinate geometry)
% ------------------------------------------------------------------
% \distnote
% \begin{align*}
% \dist{AB}{3}{1}{6}{4}{3\sqrt2}\\     % substitution, brackets, squares, sum, root
% \distn{CD}{1}{7}{4}{2}               % stops at the unsimplified root
% \end{align*}

\end{document}
